\section{Angels in Sacred Scripture and the Faith of the Church}
\label{sec:scripture-and-faith}

In the year that King Ozias died, the prophet Isaias beheld the Lord upon a high and elevated throne, His train filling the temple, and six-winged seraphim stationed over Him:
\begin{quotation}
Upon it stood the seraphims: the one had six wings, and the other had six wings: with two they covered his face, and with two they covered his feet, and with two they flew. And they cried one to another, and said: Holy, holy, holy, the Lord God of hosts, all the earth is full of his glory, And the lintels of the doors were moved at the voice of him that cried, and the house was filled with smoke. (Is~6:2--4)
\end{quotation}
Around the divine throne stands an innumerable city of living, rational spirits who gaze upon God, echo His holiness, and run upon His errands among men. From the flaming sword before Eden to the angel who checks St.~John's prostration on Patmos, they belong wholly to Christ and to the salvation of His Church.

\subsection{Ministers of Theophany and Covenant in the Old Testament}
\label{subsec:ot-theophany}

At the threshold of human exile, God ``placed before the paradise of pleasure Cherubims, and a flaming sword, turning every way, to keep the way of the tree of life'' (Gen~3:24). The road back to immortality is kept by celestial watchmen until God Himself opens it.

In the covenant with Abraham, angels walk Canaan in human guise. Sitting at his tent door in the heat of the day in the vale of Mambre, Abraham sees three men standing near him, runs to bow to the ground, washes their feet, and sets butter, milk, and a boiled calf before them; under the tree, the Lord promises that Sara will bear a son (Gen~18:1--15). Two of the three walk to Sodom at evening, where Lot presses them into his house and bakes unleavened bread before they lead him from the fire (Gen~19:1--3). Centuries later, the Epistle to the Hebrews recalls those tents: ``And hospitality do not forget: for by this some, being not aware of it, have entertained angels'' (Heb~13:2).

To Jacob sleeping on a stone at Bethel, heaven opens: ``And he saw in his sleep a ladder standing upon the earth, and the top thereof touching heaven: the angels also of God ascending and descending by it'' (Gen~28:12). Waking, he names the spot ``the house of God, and the gate of heaven'' (Gen~28:17). Returning to Canaan years later, ``the angels of God met him,'' and he called that camp Mahanaim, saying, ``These are the camps of God'' (Gen~32:1--2 [Heb~32:2--3]).

During the Exodus and conquest, angels embody God's faithful guidance. At the burning bush on Horeb (Ex~3:1--6), the Clementine Vulgate and Douay-Rheims read in verse~2, ``And the Lord appeared to him in a flame of fire out of the midst of a bush,'' while the Hebrew and the Septuagint name the Angel of the Lord in the flame before the voice identifies Himself as the God of Abraham, Isaac, and Jacob. At Sinai, God promises an angelic escort:
\begin{quotation}
Behold I will send my angel, who shall go before thee, and keep thee in thy journey, and bring thee into the place that I have prepared. Take notice of him, and hear his voice, and do not think him one to be contemned: for he will not forgive when thou hast sinned, and my name is in him. (Ex~23:20--21)
\end{quotation}
Against disobedience, the messenger bears a drawn sword. When Balaam rides to curse Israel, an angel of the Lord bars the narrow vineyard way; the ass sees him and drops beneath her rider before the Lord opens Balaam's eyes to fall flat upon the ground (Num~22:22--35). Outside Jericho, Josue meets a man with a drawn sword who answers neither as ally nor as foe, ``No: but I am prince of the host of the Lord, and now I am come,'' and bids him loose his shoes on holy ground (Jos~5:13--16 [Heb~5:13--15]).

In the Books of Kings, angels serve both judgment and mercy. After David's census, the Lord stays the destroying angel's hand beside the threshing-floor of Areuna the Jebusite (2~Kings~24:16--17 [2~Sam~24:16--17]). In the desert south of Bersabee, when Elias sleeps under a juniper tree, an angel twice wakes him with a hearth cake and water---``Arise, eat: for thou hast yet a great way to go''---and in that food's strength he walks forty days to Horeb (3~Kings~19:5--8 [1~Kgs~19:5--8]).

\subsection{The Heavenly Court, the Psalter, and the Prophetic Visions}
\label{subsec:ot-court-prophets}

In the Book of Job, ``the sons of God'' assemble before the Lord, with Satan appearing strictly under divine limit: ``only put not forth thy hand upon his person'' (Job~1:6--12). Even celestial ministers stand mutable before God, who ``in his angels... found wickedness'' (Job~4:18), yet they already praised Him when earth's foundations were laid: ``When the morning stars praised me together, and all the sons of God made a joyful melody?'' (Job~38:7).

Israel's Psalter joins human prayer to that heavenly court. Man stands just beneath the spirits: ``Thou hast made him a little less than the angels, thou hast crowned him with glory and honour'' (Ps~8:6 [Heb~8:6; Eng.~8:5]). Those who fear God walk under an angelic guard: ``The angel of the Lord shall encamp round about them that fear him: and shall deliver them'' (Ps~33:8 [Heb~34:8; Eng.~34:7]), and ``he hath given his angels charge over thee; to keep thee in all thy ways. In their hands they shall bear thee up: lest thou dash thy foot against a stone'' (Ps~90:11--12 [91:11--12]). Desert manna is ``the bread of angels'' (Ps~77:25 [78:25]). Against every idol stands the command later applied to Christ: ``Adore him, all you his angels'' (Ps~96:7 [97:7]). The psalmist summons those ``mighty in strength, and execute his word... ministers of his that do his will'' (Ps~102:20--21 [103:20--22]), blessing Him ``Who makest thy angels spirits: and thy ministers a burning fire'' (Ps~103:4 [104:4]). Psalm~148 places the angels at the head of all creation, squarely on the creaturely side of the line: ``Praise ye him, all his angels, praise ye him, all his hosts... For he spoke, and they were made: he commanded, and they were created'' (Ps~148:2, 5).

In the prophets, the heavenly court blazes forth in vision and personal names. In Isaias~6:6--8, a seraph flies from the altar with a live coal held in tongs to cleanse the prophet's lips. Later in Isaias, the taunt against Babylon---``How art thou fallen from heaven, O Lucifer, who didst rise in the morning?... I will be like the most High'' (Is~14:12--14)---stands beside Ezechiel's lament over Tyre, ``Thou a cherub stretched out, and protecting... Thou wast perfect in thy ways from the day of thy creation, until iniquity was found in thee'' (Ezech~28:14--15), passages in which the Fathers read the figure of the devil's proud fall.

By the river Chobar, Ezechiel beholds four living creatures amid fire and amber, each bearing four wings and four faces---man, lion, ox, and eagle---moving straight forward beside towering, eye-filled wheels ``as it were a wheel in the midst of a wheel,'' beneath a crystal firmament and sapphire throne (Ezech~1:4--28). In the temple vision of chapter~10, the prophet recognizes them---``I understood that they were cherubims'' (Ezech~10:20)---and watches one cherub hand burning coals from between the wheels to the angel clothed in linen (Ezech~10:2, 7).

Daniel, Tobias, and Zacharias supply the only three proper angelic names in canonical Scripture. Daniel sees the Ancient of Days enthroned in flame where ``thousands of thousands ministered to him, and ten thousand times a hundred thousand stood before him'' (Dan~7:9--10). A voice directs \textbf{Gabriel}, ``make this man to understand the vision'' (Dan~8:16), and Gabriel returns ``flying swiftly'' at the evening sacrifice to instruct him (Dan~9:21--23). In chapter~10, a messenger reports that ``the prince of the kingdom of the Persians resisted me one and twenty days: and behold \textbf{Michael}, one of the chief princes, came to help me'' (Dan~10:13), adding that none aids him against the princes of Persia and Greece ``but Michael your prince'' (Dan~10:20--21), who shall rise at the end as ``the great prince, who standeth for the children of thy people'' (Dan~12:1).

In Tobias, God hears the prayers of blind Tobit and afflicted Sara and sends ``the holy angel of the Lord, \textbf{Raphael}'' to heal both (Tob~3:24--25 [LXX~3:16--17]). Walking as ``a beautiful young man, standing girded'' (Tob~5:5 [LXX~5:4]), Raphael guides the younger Tobias, frees Sara, restores Tobit's sight, and unveils himself at parting:
\begin{quotation}
When thou didst pray with tears, and didst bury the dead... I offered thy prayer to the Lord... For I am the angel Raphael, one of the seven, who stand before the Lord... I seemed indeed to eat and to drink with you but I use an invisible meat and drink, which cannot be seen by men. (Tob~12:12, 15, 19)
\end{quotation}
Zacharias hears an angel speaking within him, listens as the angel of the Lord pleads for Jerusalem (Zach~1:9--14), and sees the angel defending Josue the high priest while Satan stands at his right hand until the Lord rebukes him (Zach~3:1--2).

\subsection{The Angels Around Christ in the New Testament}
\label{subsec:nt-christ}

In the New Testament, every angelic ministry converges on Jesus Christ. At the start of His public life, Christ claims Jacob's ladder for Himself: ``Amen, amen, I say to you, you shall see the heaven opened and the angels of God ascending and descending upon the Son of man'' (John~1:51).

At the temple altar of incense, Gabriel announces John's birth to Zachary: ``I am Gabriel, who stand before God and am sent to speak to thee and to bring thee these good tidings'' (Luke~1:11--19). Six months later he greets the Virgin Mary at Nazareth---``Hail, full of grace, the Lord is with thee: blessed art thou among women''---declaring that the Holy Ghost will overshadow her to conceive the Son of the Most High, and departing upon her assent (Luke~1:26--38). An angel reassures Joseph in sleep (Matt~1:20). On the night of the Nativity, God's brightness surrounds the shepherds as an angel proclaims ``a Saviour, who is Christ the Lord,'' joined suddenly by ``a multitude of the heavenly army, praising God and saying: Glory to God in the highest: and on earth peace to men of good will'' (Luke~2:8--15).

In the desert, Satan quotes Psalm~90:11--12 to tempt Jesus; when Christ drives him away, ``behold angels came and ministered to him'' (Matt~4:1--11). In Gethsemane, ``there appeared to him an angel from heaven, strengthening him'' in His agony (Luke~22:43), though moments later He tells Peter His Father would grant Him ``presently more than twelve legions of angels'' (Matt~26:53).

On Easter morning, an angel descends amid an earthquake with countenance as lightning and raiment as snow, rolls back the stone, and announces, ``He is not here. For he is risen, as he said'' (Matt~28:2--6). St.~Luke records two men in shining apparel asking, ``Why seek you the living with the dead?'' (Luke~24:4--8), and Mary Magdalene sees ``two angels in white, sitting, one at the head, and one at the feet, where the body of Jesus had been laid'' (John~20:12), like the cherubim over the ark's mercy seat. In Acts, angels open prison doors for the apostles (Acts~5:19) and loose Peter's chains by night; hearing Peter's voice at the gate, the praying disciples first say, ``It is his angel'' (Acts~12:7--15). Against the Sadducees, who deny ``resurrection, neither angel, nor spirit'' (Acts~23:8), the apostolic faith confesses both.

Christ Himself teaches personal guardianship, heavenly joy over penance, and angelic judgment:
\begin{quotation}
See that you despise not one of these little ones: for I say to you, that their angels in heaven always see the face of my Father who is in heaven. (Matt~18:10)
\end{quotation}
There is ``joy before the angels of God upon one sinner doing penance'' (Luke~15:10), and Lazarus dies to be ``carried by the angels into Abraham's bosom'' (Luke~16:22). Risen saints ``shall be as the angels of God in heaven'' (Matt~22:30), ``equal to the angels'' in deathlessness (Luke~20:36). Yet even angels know not the final day and hour (Matt~24:36). When the Son of Man comes in majesty with all the angels (Matt~25:31), ``with the voice of an archangel and with the trumpet of God'' (1~Thess~4:15 [Vulgate/Greek~4:16]) and ``with the angels of his power'' (2~Thess~1:7), the angels reap and separate wicked from just (Matt~13:39--49), and the damned go into ``everlasting fire, which was prepared for the devil and his angels'' (Matt~25:41).

The Epistles establish Christ's supremacy over every spirit. St.~Paul charges Timothy before ``God and Christ Jesus and the elect angels'' (1~Tim~5:21), calls the apostles ``a spectacle to the world and to angels and to men'' (1~Cor~4:9), notes reverence in assembly ``because of the angels'' (1~Cor~11:10), invokes ``the tongues of men and of angels'' (1~Cor~13:1), and asks whether believers know ``that we shall judge angels'' (1~Cor~6:3). Though the Old Law was ``ordained by angels'' (Gal~3:19; Acts~7:53), even an angel preaching another gospel is anathema (Gal~1:8), and neither ``angels, nor principalities, nor powers'' can sever us from Christ's love (Rom~8:38--39).

Ephesians and Colossians supply the Pauline names later gathered into nine choirs while setting Christ above them all. The Father seated the risen Christ ``Above all principality and power and virtue and dominion and every name that is named'' (Eph~1:20--21). Through the Church, God's manifold wisdom is made known ``to the principalities and powers in heavenly places'' (Eph~3:10), even while Christians wrestle ``against principalities and powers, against the rulers of the world of this darkness, against the spirits of wickedness in the high places'' (Eph~6:12). Against a false ``religion of angels'' detached from the Head (Col~2:18--19), Paul declares:
\begin{quotation}
For in him were all things created in heaven and on earth, visible and invisible, whether thrones, or dominations, or principalities, or powers. All things were created by him and in him. And he is before all: and by him all things consist. (Col~1:16--17)
\end{quotation}
Hebrews sharpens that divide: Christ inherits a name infinitely above the angels, who are commanded, ``And let all the angels of God adore him'' (Heb~1:4--6, citing Ps~96:7), whereas angels are mutable servants---``spirits'' and ``a flame of fire'' (Heb~1:7, citing Ps~103:4)---``all ministering spirits, sent to minister for them who shall receive the inheritance of salvation'' (Heb~1:14). Jesus was made ``a little lower than the angels, for the suffering of death, crowned with glory and honour'' (Heb~2:5--9), opening to Christians ``the heavenly Jerusalem, and... the company of many thousands of angels'' (Heb~12:22--24).

Meanwhile, the devil was ``a murderer from the beginning,'' who ``stood not in the truth'' (John~8:44). God ``spared not the angels that sinned'' (2~Pet~2:4); ``the angels who kept not their principality'' are reserved in darkness unto judgment (Jude~1:6), while ``Michael the archangel, disputing with the devil'' over Moses' body says only, ``The Lord command thee'' (Jude~1:9), and Jude cites Enoch's prophecy of the Lord coming ``with thousands of his saints'' (Jude~1:14--15).

In the Apocalypse, Christ holds ``the angels of the seven churches'' (Apoc~1:20). Around the throne, four six-winged living creatures full of eyes chant unceasingly ``Holy, Holy, Holy, Lord God Almighty'' (Apoc~4:6--8), joined by ``thousands of thousands'' of angels (Apoc~5:11). Four angels restrain the winds while another seals God's servants on their foreheads, and all the angels fall prostrate in adoration (Apoc~7:1--12). Seven angels stand with seven trumpets, while another lifts ``the smoke of the incense of the prayers of the saints'' from a golden censer (Apoc~8:2--4). In heaven, ``Michael and his angels fought with the dragon, and the dragon fought, and his angels,'' hurling down the accuser, conquered ``by the blood of the Lamb and by the word of the testimony'' (Apoc~12:7--12). Twice St.~John falls at his angelic guide's feet, and twice the angel checks him: ``See thou do it not. I am thy fellow servant and of thy brethren who have the testimony of Jesus. Adore God'' (Apoc~19:10; cf.~22:8--9).

\subsection{The Confession of the Church}
\label{subsec:church-confession}

In the Nicene-Constantinopolitan Creed (325/381), the Church confesses God the Father Almighty as Maker of heaven and earth, \latin{visibilium omnium et invisibilium}, ``of all things visible and invisible.'' The angels are confessed at once as real beings and as creatures made out of nothing by the Triune God, utterly distinct from the divine essence.

Early synods guarded this faith against idolatry and speculative distortion. The Council of Laodicea (c.~343--381), in canon~35, anathematized Christians who forsake the Church of God to invoke angels in a hidden idolatry, abandoning our Lord Jesus Christ. Against Origenist cosmology, the decrees of Constantinople (543) and the Fifth Ecumenical Council (Constantinople~II, 553) rejected the fable that pre-existent rational minds fell through satiety into angelic, human, and demonic ranks; condemned the claim that the Word became an angel among angels or would be crucified for demons; and anathematized the teaching that the punishment of demons and impious men is temporary and will end in a universal restoration (\gk{apokatastasis}).

Against Cathar dualism, the Fourth Lateran Council (1215) promulgated the Church's classic solemn definition in Constitution~1, \emph{Firmiter credimus} (DS~800--801), defining that the one true Triune God is:
\begin{quotation}
\latin{unum universorum principium: creator omnium visibilium et invisibilium, spiritualium et corporalium: qui sua omnipotenti virtute simul ab initio temporis utramque de nihilo condidit creaturam, spiritualem et corporalem, angelicam videlicet et mundanam: ac deinde humanam, quasi communem ex spiritu et corpore constitutam. Diabolus enim et alii daemones a Deo quidem natura creati sunt boni, sed ipsi per se facti sunt mali. Homo vero diaboli suggestione peccavit.} (Lateran~IV, Const.~1, DS~800)
\end{quotation}
That is: ``one sole principle of all things, Creator of all things visible and invisible, spiritual and corporeal; who by His almighty power together from the beginning of time made out of nothing both orders of creation, the spiritual and the corporeal, namely the angelic and the worldly; and afterward the human, constituted as it were in common of spirit and body. For the devil and the other demons were indeed created naturally good by God, yet became evil by their own doing. Man, however, sinned at the suggestion of the devil.'' At the end of time, the council adds (DS~801), the reprobate receive \latin{cum diabolo poenam perpetuam}, ``with the devil perpetual punishment,'' and the elect \latin{cum Christo gloriam sempiternam}, ``with Christ everlasting glory.'' Four points stand solemnly defined: the existence of spiritual creatures; their creation \latin{ex nihilo} together with the material cosmos at the beginning of time; the original natural goodness of the devil and every demon; and the voluntary, irreversible origin of demonic malice.

The Council of Florence expanded both sides of this confession. In \emph{Laetentur caeli} (1439, DS~1305), it defined that purified souls in heaven behold clearly the Triune God as He is, according to diversity of merits. In \emph{Cantate Domino} (1442, DS~1333, 1336), it reaffirmed against Manichaean dualism that the one true God created all spiritual and corporeal creatures out of nothing---good because made by the Highest Good, mutable because made from nothing---and declared that no nature of evil exists, since every nature insofar as it is nature is good.

The Council of Trent confessed humanity's captivity under ``the empire of death, that is to say the devil'' (Session~5, Original Sin, canon~1) and the ongoing combat of the justified ``with the flesh, with the world, with the devil'' (Session~6, Justification, ch.~13), while upholding in Session~25 the relative veneration of sacred images and ordering superstition purged. The \emph{Roman Catechism} (1566) expounded this faith for pastors: God ``created from nothing spiritual nature, and angels innumerable to serve and minister to him'' (Part~I, Creed, art.~1); the faithful lawfully venerate and invoke angels without infringing the adoration (\gk{latria}) owed to God alone, since ``they always see the face of God, and are constituted by him the willing advocates of our salvation'' (Part~III, First Commandment); ``Angels are commissioned by Divine Providence to guard the human race, and be present with every man to protect him from injury'' (Part~IV, Preface to the Lord's Prayer); and their joyful obedience models the petition ``Thy will be done on earth as it is in heaven.''

In 1870, the First Vatican Council in \emph{Dei Filius} (ch.~1) repeated Lateran~IV's formula verbatim, defining that God freely and to manifest His perfection \latin{simul ab initio temporis utramque de nihilo condidit creaturam, spiritualem et corporalem, angelicam videlicet et mundanam}, ``together from the beginning of time made out of nothing both orders of creation, the spiritual and the corporeal, namely the angelic and the worldly.'' Canons~I.1--5 anathematize anyone who denies the Creator of visible and invisible things (I.1), affirms \latin{praeter materiam nihil esse}, ``that nothing exists besides matter'' (I.2), identifies God's substance with creatures (I.3), claims finite spiritual things emanated from the divine substance (I.4), or denies that all things, \latin{et spirituales et materiales}, ``both spiritual and material,'' were produced freely out of nothing (I.5).

The twentieth-century Magisterium reaffirmed this personal realism with equal clarity. Pope Pius~XII in \emph{Humani generis} (1950, \S 26) rejected the opinion calling into question whether the angels are personal beings. The Second Vatican Council in \emph{Lumen gentium} (1964, \S\S 53, 56, 59, 66) confessed Mary's pre-eminence above all creatures in heaven and on earth as Mother of God and Queen of the universe, while preserving the distinction between divine adoration and the veneration of Mary, the saints, and the angels. The Congregation for the Doctrine of the Faith in \emph{Christian Faith and Demonology} (1975) reaffirmed Lateran~IV's teaching on Satan and the demons as real, personal creatures created good by God who became evil by their own choice. And the \emph{Catechism of the Catholic Church} (\S\S 325--336, 391--395) summarizes the whole confession: the existence of spiritual, non-corporeal, personal, and immortal beings is a truth of faith; the angels surround Christ, join the Church's liturgy, and guard every believer; and Satan and the other demons are fallen angels whose choice against God is irrevocable, yet whose power remains finite and governed by divine providence.
